\chap{Введение}

Цели данной практической работы:
\begin{enumerate}
	\item изучение принципов управления пакетами в операционных системах семейства GNU/Linux;
	\item освоение практических навыков работы с пакетными менеджерами APT и dpkg;
	\item изучение различных типов пакетов (deb, метапакеты, snap, виртуальные);
	\item получение опыта создания собственного .deb-пакета, настройки локального репозитория и установки пакетов из него.
\end{enumerate}

\chap{Выполнение}

\sect{Задание 1. Идентификация типов пакетов}

\begin{table}[!htbp]
	\caption{Типы пакетов}
	\begin{tabularx}{\linewidth}{|X|X|X|X|X|}
		\hline
		\textbf{Пакет} & \textbf{Тип} & \textbf{Размер} & \textbf{Зависимости} & С\textbf{одержит файлы} \\\hline
		nano &&&& \\\hline
		ubuntu-minimal &&&& \\\hline
		firefox &&&& \\\hline
		core22 &&&& \\\hline
		chromium &&&& \\\hline
	\end{tabularx}
\end{table}

Вопросы по заданию:
\begin{enumerate}
	\item \textit{Какие типы пакетов вы обнаружили? Перечислите их.}
	\item \textit{Чем отличается вывод apt show для обычного пакета и метапакета?}
	\item \textit{Почему для firefox команда dpkg -L показывает мало файлов, но браузер работает?}
	\item \textit{Что такое виртуальный пакет? Приведите пример из задания.}
	\item \textit{Как определить, что пакет является snap-пакетом, используя данные команды?}
\end{enumerate}

\sect{Задание 2. Анализ метапакетов в системе}

Вопросы по заданию:
\begin{enumerate}
	\item \textit{Какие метапакеты вы обнаружили в системе?}
	\item \textit{Какие три типа зависимостей (Depends/Recommends/Suggests) присутствуют у ubuntu-standard? Приведите по одному примеру каждого типа.}
	\item \textit{Сколько файлов содержит метапакет ubuntu-standard? Почему их так мало?}
	\item \textit{Был ли `ubuntu-standard` уже установлен в системе? Если да, что это говорит о вашей системе?}
	\item \textit{Какой пакет вы установили дополнительно? Какие у него были зависимости?}
	\item \textit{Что произошло при удалении дополнительного пакета? Удалились ли его зависимости?}
	\item \textit{Какие практические преимущества дает использование метапакетов при развертывании одинаковых систем на нескольких компьютерах?}
\end{enumerate}

\sect{Задание 3. Создание простейшего .deb пакета}

Вопросы по заданию:
\begin{enumerate}
	\item \textit{Какие файлы обязательно должны быть в .deb пакете?}
	\item \textit{Почему важно указывать архитектуру (Architecture) в control-файле?}
	\item \textit{Как удалить созданный пакет из системы?}
\end{enumerate}

\sect{Задание 4. Настройка локального репозитория}

Вопросы по заданию:
\begin{enumerate}
	\item \textit{Для чего нужна утилита dpkg-scanpackages и что она создаёт?}
	\item \textit{Что представляет собой файл Packages.gz и какую роль он играет в работе APT?}
	\item \textit{Почему в строке репозитория необходимо указывать параметр [trusted=yes]? Что произойдёт без него?}
	\item \textit{Как проверить, что локальный репозиторий успешно добавлен и пакет в нём доступен для установки?}
	\item \textit{Какие преимущества даёт использование локального репозитория по сравнению с установкой пакета напрямую через dpkg -i?}
	\item \textit{В каких ситуациях системному администратору может потребоваться создавать локальный репозиторий?}
	\item \textit{Как удалить локальный репозиторий из системы, чтобы он больше не использовался?}
\end{enumerate}

\chap{Вывод}

В ходе практической работы мы:
\begin{itemize}
	\item изучили принципы управления пакетами в операционных системах семейства GNU/Linux;
	\item освоили практические навыки работы с пакетными менеджерами APT и dpkg;
	\item изучили различные типы пакетов;
	\item получили опыт создания собственного .deb-пакета, настройки локального репозитория и установки пакетов из него.
\end{itemize}
